Table~\ref{tab:main} gives the main comparison (seeds 0--4, test split), Table~\ref{tab:paired} the paired differences in average accuracy between ER with $M{=}100$ and the other systems, and Fig.~\ref{fig:curves} the mean accuracy on the tasks seen so far after each stage. Means quoted below are from Table~\ref{tab:main}. Differences and $p$-values are those recorded by \texttt{rh compare} with ER ($M{=}100$) as the reference system (Table~\ref{tab:paired}): a difference is ER with $M{=}100$ minus the other system, and $p$ is that of the paired $t$-test.

\begin{table*}[t]\centering\caption{Main results on the test split, mean $\pm$ std over 5 seeds (seeds 0--4). EWC uses the validation-selected $\lambda$ per scenario; experience replay is shown at $M{=}20$ and $M{=}100$. Bold: best printed mean in the scenario (all tied systems are bold); underline: runner-up, shown only when the best is not tied. Joint training is the upper-bound reference, and its backward transfer compares networks retrained from scratch.}\label{tab:main}
\begin{tabular}{llccccc}
\toprule
Method & Scenario & Avg. accuracy $\uparrow$ & Forgetting $\downarrow$ & Backward transfer $\uparrow$ & Last-task acc. $\uparrow$ & Seeds \\
\midrule
Fine-tuning & perm\_dil & 0.498 $\pm$ 0.021 & 0.584 $\pm$ 0.023 & $-$0.584 $\pm$ 0.023 & 0.973 $\pm$ 0.009 & 5 \\
EWC & perm\_dil & 0.695 $\pm$ 0.019 & 0.303 $\pm$ 0.024 & $-$0.303 $\pm$ 0.024 & 0.917 $\pm$ 0.017 & 5 \\
Experience replay (M=20) & perm\_dil & 0.670 $\pm$ 0.042 & 0.370 $\pm$ 0.052 & $-$0.370 $\pm$ 0.052 & \underline{0.975 $\pm$ 0.009} & 5 \\
Experience replay (M=100) & perm\_dil & \underline{0.847 $\pm$ 0.008} & \underline{0.148 $\pm$ 0.009} & \underline{$-$0.148 $\pm$ 0.009} & \textbf{0.976 $\pm$ 0.011} & 5 \\
Joint training & perm\_dil & \textbf{0.954 $\pm$ 0.010} & \textbf{0.010 $\pm$ 0.006} & \textbf{$-$0.004 $\pm$ 0.005} & 0.967 $\pm$ 0.009 & 5 \\
\midrule
Fine-tuning & split\_cil & 0.214 $\pm$ 0.031 & 0.973 $\pm$ 0.039 & $-$0.973 $\pm$ 0.039 & \textbf{0.997 $\pm$ 0.006} & 5 \\
EWC & split\_cil & 0.246 $\pm$ 0.050 & 0.568 $\pm$ 0.136 & $-$0.568 $\pm$ 0.136 & 0.118 $\pm$ 0.263 & 5 \\
Experience replay (M=20) & split\_cil & 0.728 $\pm$ 0.051 & 0.327 $\pm$ 0.062 & $-$0.327 $\pm$ 0.062 & \underline{0.992 $\pm$ 0.012} & 5 \\
Experience replay (M=100) & split\_cil & \underline{0.896 $\pm$ 0.021} & \underline{0.112 $\pm$ 0.032} & \underline{$-$0.112 $\pm$ 0.032} & 0.986 $\pm$ 0.010 & 5 \\
Joint training & split\_cil & \textbf{0.957 $\pm$ 0.015} & \textbf{0.028 $\pm$ 0.018} & \textbf{$-$0.026 $\pm$ 0.020} & 0.964 $\pm$ 0.018 & 5 \\
\midrule
Fine-tuning & split\_til & 0.946 $\pm$ 0.025 & 0.059 $\pm$ 0.028 & $-$0.058 $\pm$ 0.029 & \textbf{0.997 $\pm$ 0.006} & 5 \\
EWC & split\_til & 0.973 $\pm$ 0.015 & \underline{0.006 $\pm$ 0.005} & \underline{$-$0.005 $\pm$ 0.004} & 0.945 $\pm$ 0.071 & 5 \\
Experience replay (M=20) & split\_til & 0.973 $\pm$ 0.011 & 0.026 $\pm$ 0.014 & $-$0.025 $\pm$ 0.014 & \textbf{0.997 $\pm$ 0.006} & 5 \\
Experience replay (M=100) & split\_til & \underline{0.984 $\pm$ 0.006} & 0.012 $\pm$ 0.009 & $-$0.010 $\pm$ 0.010 & \textbf{0.997 $\pm$ 0.006} & 5 \\
Joint training & split\_til & \textbf{0.990 $\pm$ 0.009} & \textbf{0.005 $\pm$ 0.005} & \textbf{$-$0.003 $\pm$ 0.005} & \textbf{0.997 $\pm$ 0.006} & 5 \\
\bottomrule
\end{tabular}

\end{table*}

\begin{table}[t]\centering\caption{Paired differences in average accuracy over seeds 0--4 (main runs), from \texttt{rh compare} with ER ($M{=}100$) as reference: mean difference (ER with $M{=}100$ minus the other system) and paired $t$-test $p$ (uncorrected). FT: fine-tuning.}\label{tab:paired}
\begin{tabular}{llcc}
\toprule
Scenario & Other system & Diff. & $p$ \\
\midrule
perm\_dil & EWC & \rhval{cmp/main/ewc/perm_dil/average_accuracy/delta:3} & \rhval{cmp/main/ewc/perm_dil/average_accuracy/paired_p} \\
perm\_dil & Joint & \rhval{cmp/main/joint-training/perm_dil/average_accuracy/delta:3} & \rhval{cmp/main/joint-training/perm_dil/average_accuracy/paired_p} \\
\midrule
split\_cil & FT & \rhval{cmp/main/fine-tuning/split_cil/average_accuracy/delta:3} & \rhval{cmp/main/fine-tuning/split_cil/average_accuracy/paired_p} \\
split\_cil & EWC & \rhval{cmp/main/ewc/split_cil/average_accuracy/delta:3} & \rhval{cmp/main/ewc/split_cil/average_accuracy/paired_p} \\
split\_cil & Joint & \rhval{cmp/main/joint-training/split_cil/average_accuracy/delta:3} & \rhval{cmp/main/joint-training/split_cil/average_accuracy/paired_p} \\
\midrule
split\_til & FT & \rhval{cmp/main/fine-tuning/split_til/average_accuracy/delta:3} & \rhval{cmp/main/fine-tuning/split_til/average_accuracy/paired_p} \\
split\_til & EWC & \rhval{cmp/main/ewc/split_til/average_accuracy/delta:3} & \rhval{cmp/main/ewc/split_til/average_accuracy/paired_p} \\
split\_til & Joint & \rhval{cmp/main/joint-training/split_til/average_accuracy/delta:3} & \rhval{cmp/main/joint-training/split_til/average_accuracy/paired_p} \\
\bottomrule
\end{tabular}
\end{table}

\begin{figure*}[t]\centering\includegraphics[width=0.92\textwidth]{curves_main.pdf}
\caption{Mean accuracy on the tasks seen so far after each stage (main runs, mean $\pm$ std over 5 seeds).}\label{fig:curves}\end{figure*}

\textbf{H1 (permuted digits, ER vs EWC): supported.} ER with $M{=}100$ reaches \rhval{main/experience-replay-m-100/perm_dil/average_accuracy/mean:3} average accuracy against \rhval{main/ewc/perm_dil/average_accuracy/mean:3} for EWC (difference \rhval{cmp/main/ewc/perm_dil/average_accuracy/delta:3}, $p{=}$\rhval{cmp/main/ewc/perm_dil/average_accuracy/paired_p}) and has lower forgetting (\rhval{main/experience-replay-m-100/perm_dil/forgetting/mean:3} vs \rhval{main/ewc/perm_dil/forgetting/mean:3}). EWC also has the lowest last-task accuracy of the five systems, \rhval{main/ewc/perm_dil/last_task_accuracy/mean:3} against \rhval{main/experience-replay-m-100/perm_dil/last_task_accuracy/mean:3} for ER and \rhval{main/fine-tuning/perm_dil/last_task_accuracy/mean:3} for fine-tuning. The claim does not extend to small buffers: ER with $M{=}20$ (\rhval{main/experience-replay-m-20/perm_dil/average_accuracy/mean:3}) is below EWC (\rhval{main/ewc/perm_dil/average_accuracy/mean:3}), a gap smaller than the seed standard deviation of ER with $M{=}20$ (\rhval{main/experience-replay-m-20/perm_dil/average_accuracy/std:3}), so five seeds do not resolve it. Joint training reaches \rhval{main/joint-training/perm_dil/average_accuracy/mean:3}; the difference of ER with $M{=}100$ to it is \rhval{cmp/main/joint-training/perm_dil/average_accuracy/delta:3} ($p{=}$\rhval{cmp/main/joint-training/perm_dil/average_accuracy/paired_p}), so a 100-sample buffer does not close the gap to joint training.

\textbf{H2 (class-incremental split digits): partly supported, seed-dependent.} Fine-tuning gets \rhval{main/fine-tuning/split_cil/average_accuracy/mean:3} with forgetting \rhval{main/fine-tuning/split_cil/forgetting/mean:3} and last-task accuracy \rhval{main/fine-tuning/split_cil/last_task_accuracy/mean:3}, the pattern expected of a network that predicts only the two most recent classes. EWC with $\lambda{=}10^3$ reaches \rhval{main/ewc/split_cil/average_accuracy/mean:3}, a gain below the registered 0.05 threshold but within noise (the seed standard deviation of EWC is \rhval{main/ewc/split_cil/average_accuracy/std:3}). The EWC half of H2 therefore holds only by the registered mean criterion on the main seeds: with the identical configuration ($\lambda{=}10^3$) on the sweep seeds 10--14, EWC reaches \rhval{sweep_lambda/ewc@lam=1000/split_cil/average_accuracy/mean:3} against \rhval{sweep_lambda/ewc@lam=0/split_cil/average_accuracy/mean:3} for fine-tuning (Table~\ref{tab:sweep}), a gain well above 0.05. We do not reconcile the two seed sets beyond noting that the EWC result on this scenario is highly seed-dependent. The ER half holds: ER with $M{=}100$ reaches \rhval{main/experience-replay-m-100/split_cil/average_accuracy/mean:3} (gain \rhval{cmp/main/fine-tuning/split_cil/average_accuracy/delta:3} over fine-tuning, $p{=}$\rhval{cmp/main/fine-tuning/split_cil/average_accuracy/paired_p}, and \rhval{cmp/main/ewc/split_cil/average_accuracy/delta:3} over EWC, $p{=}$\rhval{cmp/main/ewc/split_cil/average_accuracy/paired_p}), and ER with $M{=}20$ reaches \rhval{main/experience-replay-m-20/split_cil/average_accuracy/mean:3}. This is consistent with class-incremental learning being the hardest scenario in \cite{ven2019three}.

\textbf{Failure case: EWC on split\_cil.} EWC does not merely fail to gain here. Its last-task accuracy is \rhval{main/ewc/split_cil/last_task_accuracy/mean:3} $\pm$ \rhval{main/ewc/split_cil/last_task_accuracy/std:3}: in the run registry it is exactly zero in four of the five seeds, that is, no test image of the fifth class pair is classified correctly after training on it, and non-zero in the remaining seed. At this $\lambda$ the network therefore keeps part of the earlier tasks (forgetting \rhval{main/ewc/split_cil/forgetting/mean:3} against \rhval{main/fine-tuning/split_cil/forgetting/mean:3} for fine-tuning) and does not acquire the last one, which leaves the average near the fine-tuning level. We did not run an experiment that isolates the cause; that the summed penalties of four earlier tasks are too stiff for the single shared output layer may explain it, but this remains a hypothesis. The sweep (Table~\ref{tab:sweep}) shows no $\lambda$ above \rhval{sweep_lambda/ewc@lam=1000/split_cil/average_accuracy/mean:3} average accuracy on this scenario.

\textbf{H3 (task-incremental split digits): supported, but the scenario is nearly saturated.} Fine-tuning reaches \rhval{main/fine-tuning/split_til/average_accuracy/mean:3} and joint training \rhval{main/joint-training/split_til/average_accuracy/mean:3}, a gap inside the registered 0.10. ER ($M{=}100$) \rhval{main/experience-replay-m-100/split_til/average_accuracy/mean:3}, ER ($M{=}20$) \rhval{main/experience-replay-m-20/split_til/average_accuracy/mean:3} and EWC \rhval{main/ewc/split_til/average_accuracy/mean:3} lie between them. ER with $M{=}100$ is not separated from fine-tuning at the 0.05 level ($p{=}$\rhval{cmp/main/fine-tuning/split_til/average_accuracy/paired_p}), from EWC ($p{=}$\rhval{cmp/main/ewc/split_til/average_accuracy/paired_p}) or from joint training ($p{=}$\rhval{cmp/main/joint-training/split_til/average_accuracy/paired_p}). We therefore cannot rank the three continual methods here. EWC's last-task accuracy (\rhval{main/ewc/split_til/last_task_accuracy/mean:3} $\pm$ \rhval{main/ewc/split_til/last_task_accuracy/std:3}) is below the \rhval{main/experience-replay-m-100/split_til/last_task_accuracy/mean:3} that the other four systems share.

\textbf{H5 (joint upper bound): supported in the means, not significantly on split\_til.} Joint training has the highest average accuracy in every scenario (\rhval{main/joint-training/perm_dil/average_accuracy/mean:3}, \rhval{main/joint-training/split_cil/average_accuracy/mean:3}, \rhval{main/joint-training/split_til/average_accuracy/mean:3}). The difference of ER with $M{=}100$ to it is \rhval{cmp/main/joint-training/perm_dil/average_accuracy/delta:3} on perm\_dil ($p{=}$\rhval{cmp/main/joint-training/perm_dil/average_accuracy/paired_p}) and \rhval{cmp/main/joint-training/split_cil/average_accuracy/delta:3} on split\_cil ($p{=}$\rhval{cmp/main/joint-training/split_cil/average_accuracy/paired_p}), and \rhval{cmp/main/joint-training/split_til/average_accuracy/delta:3} on split\_til with $p{=}$\rhval{cmp/main/joint-training/split_til/average_accuracy/paired_p}. The joint-training baseline retrains with more optimisation steps than the sequential methods, so it is an upper bound in data, not in compute.

\textbf{Forgetting versus accuracy.} Low forgetting did not imply high accuracy. On split\_til EWC has the second-lowest forgetting (\rhval{main/ewc/split_til/forgetting/mean:3}) together with the lowest last-task accuracy, and on split\_cil it has lower forgetting than fine-tuning with almost the same average accuracy. The sweep in Section~\ref{sec:ablations} shows the extreme case, forgetting near zero at the largest $\lambda$ with low average accuracy. We therefore read forgetting together with average and last-task accuracy.

\begin{figure*}[t]\centering\includegraphics[width=0.92\textwidth]{sweeps_buffer_lambda.pdf}
\caption{Final average accuracy versus buffer size $M$ (top) and EWC strength $\lambda$ (bottom), mean $\pm$ std over seeds 10--14. Dashed line: joint training mean from the main runs (different seeds).}\label{fig:sweep}\end{figure*}
